\section{Restriction and induction functors for $q$-web categories}
\label{sec:web}

In this section, we recall from \cite{SSW25} the affine $q$-web category $\qAW$ and its cyclotomic quotients $\qW_\bfu$, which are a quantum version of the affine and cyclotomic web categories introduced in \cite{SW25web}. We determine the center of $\qAW$. We formulate the restriction functors on $\qAW$ and $\qW_\bfu$ as well as the induction functors on $\qW_\bfu$. 


\subsection{Definition of $\qAW$}

Let $\kk =\C[q,q^{-1}]$.
For $n, a \in \N$, the $ q$-numbers and $q$-binomial coefficients are
\begin{align}
\label{qbinom}
    [n]_q =\frac{q^n -q^{-n}}{q-q^{-1}},
    \qquad
    \qbinom{n}{a}_q =\frac{[n]_q[n-1]_q\cdots [n-a+1]_q}{[a]_q!}.
\end{align}

\begin{definition}
  \label{def:qAWeb}
The affine $q$-web category $\qAW$ is the strict $\kk$-linear monoidal category generated by objects $a\in \mathbb Z_{\ge1}$. The object $a$ and its identity morphism  will be drawn as a vertical strand labeled by $a$,
$\stra.$
The generating morphisms are the merges, splits, (thick) positive crossings and thick dots depicted as
\begin{align}
\label{merge+split+crossing}
\merge 
&:(a,b) \rightarrow (a+b),&
\splits
&:(a+b)\rightarrow (a,b),&
\crossingpos
&:(a,b) \rightarrow (b,a), &
\wkdotaa , 
\end{align}
for $a,b \in \Z_{\ge 1}$, where $\crossingpos$ is invertible with inverse being the negative crossing denoted $\crossingneg := \crossingpos^{-1}$ and $\xdota$ is invertible with inverse denoted by $\zdota$. 
These generating morphisms are subject to the following relations \eqref{webassoc}--\eqref{dotmovesplits+merge}, for $a,b,c,d \in \Z_{\ge 1}$ with $d-a=c-b$:
\begin{align}
\label{webassoc}
\begin{tikzpicture}[baseline = 0,color=\clr]
	\draw[-,thick] (0.35,-.3) to (0.08,0.14);
	\draw[-,thick] (0.1,-.3) to (-0.04,-0.06);
	\draw[-,line width=1pt] (0.085,.14) to (-0.035,-0.06);
	\draw[-,thick] (-0.2,-.3) to (0.07,0.14);
	\draw[-,line width=1.5pt] (0.08,.45) to (0.08,.1);
        \node at (0.45,-.41) {$\scriptstyle c$};
        \node at (0.07,-.4) {$\scriptstyle b$};
        \node at (-0.28,-.41) {$\scriptstyle a$};
\end{tikzpicture}
&=
\begin{tikzpicture}[baseline = 0, color=\clr]
	\draw[-,thick] (0.36,-.3) to (0.09,0.14);
	\draw[-,thick] (0.06,-.3) to (0.2,-.05);
	\draw[-,line width=.8pt] (0.07,.14) to (0.19,-.06);
	\draw[-,thick] (-0.19,-.3) to (0.08,0.14);
	\draw[-,line width=1.5pt] (0.08,.45) to (0.08,.1);
        \node at (0.45,-.41) {$\scriptstyle c$};
        \node at (0.07,-.4) {$\scriptstyle b$};
        \node at (-0.28,-.41) {$\scriptstyle a$};
\end{tikzpicture}\:,
\qquad
\begin{tikzpicture}[baseline = -1mm, color=\clr]
	\draw[-,thick] (0.35,.3) to (0.08,-0.14);
	\draw[-,thick] (0.1,.3) to (-0.04,0.06);
	\draw[-,line width=1pt] (0.085,-.14) to (-0.035,0.06);
	\draw[-,thick] (-0.2,.3) to (0.07,-0.14);
	\draw[-,line width=1.5pt] (0.08,-.45) to (0.08,-.1);
        \node at (0.45,.4) {$\scriptstyle c$};
        \node at (0.07,.42) {$\scriptstyle b$};
        \node at (-0.28,.4) {$\scriptstyle a$};
\end{tikzpicture}
=\begin{tikzpicture}[baseline = -1mm, color=\clr]
	\draw[-,thick] (0.36,.3) to (0.09,-0.14);
	\draw[-,thick] (0.06,.3) to (0.2,.05);
	\draw[-,line width=1pt] (0.07,-.14) to (0.19,.06);
	\draw[-,thick] (-0.19,.3) to (0.08,-0.14);
	\draw[-,line width=1.5pt] (0.08,-.45) to (0.08,-.1);
        \node at (0.45,.4) {$\scriptstyle c$};
        \node at (0.07,.42) {$\scriptstyle b$};
        \node at (-0.28,.4) {$\scriptstyle a$};
\end{tikzpicture}\:,
\\
\label{mergesplit}
\begin{tikzpicture}[baseline = 7.5pt,scale=.8, color=\clr]
	\draw[-,line width=1pt] (0,0) to (.28,.3) to (.28,.7) to (0,1);
	\draw[-,line width=1pt] (.6,0) to (.31,.3) to (.31,.7) to (.6,1);
        \node at (0,1.13) {$\scriptstyle b$};
        \node at (0.63,1.13) {$\scriptstyle d$};
        \node at (0,-.1) {$\scriptstyle a$};
        \node at (0.63,-.1) {$\scriptstyle c$};
\end{tikzpicture}
&=
\sum_{\substack{0 \leq s \leq \min(a,b)\\0 \leq t \leq \min(c,d)\\t-s=d-a}}
q^{st}  
\begin{tikzpicture}[baseline = 7.5pt,scale=.8, color=\clr]
\draw[-,thick] (0.58,0) to (0.58,.2) to (.02,.8) to (.02,1);
	\draw[-,line width=4pt,white] (0.02,0) to (0.02,.2) to (.58,.8) to (.58,1);
	\draw[-,thick] (0.02,0) to (0.02,.2) to (.58,.8) to (.58,1);
	\draw[-,thick] (0,0) to (0,1);
	\draw[-,thick] (0.6,0) to (0.6,1);
        \node at (0,1.13) {$\scriptstyle b$};
        \node at (0.6,1.13) {$\scriptstyle d$};
        \node at (0,-.1) {$\scriptstyle a$};
        \node at (0.6,-.1) {$\scriptstyle c$};
        \node at (-0.1,.5) {$\scriptstyle s$};
        \node at (0.75,.5) {$\scriptstyle t$};
\end{tikzpicture},
\\
  \label{splitbinomial}
\begin{tikzpicture}[baseline = -1mm,scale=.8,color=\clr]
\draw[-,line width=1.5pt] (0.08,-.8) to (0.08,-.5);
\draw[-,line width=1.5pt] (0.08,.3) to (0.08,.6);
\draw[-,thick] (0.1,-.51) to [out=45,in=-45] (0.1,.31);
\draw[-,thick] (0.06,-.51) to [out=135,in=-135] (0.06,.31);
\node at (-.33,-.05) {$\scriptstyle a$};
\node at (.45,-.05) {$\scriptstyle b$};
\end{tikzpicture}
&= 
\qbinomsm{a+b}{a}_q\:
\begin{tikzpicture}[baseline = -1mm,scale=.6,color=\clr]
	\draw[-,line width=1.5pt] (0.08,-.8) to (0.08,.6);
        \node at (.08,-1) {$\scriptstyle a+b$};
\end{tikzpicture}, 
\\
\label{dotmovecrossing}
\begin{tikzpicture}[baseline=-1mm, scale=.8, color=\clr]
	\draw[-,line width=1pt] (0.3,-.3) to (-.3,.4);
	\draw[-,line width=4pt,white] (-0.3,-.3) to (.3,.4);
	\draw[-,line width=1pt] (-0.3,-.3) to (.3,.4);
        \node at (-0.3,-.4) {$\scriptstyle a$};
        \node at (0.3,-.4) {$\scriptstyle b$};
     \draw(-0.15,-0.12)   \bdot;
\end{tikzpicture}
 & =
\begin{tikzpicture}[baseline=-1mm, scale=.8, color=\clr]
 \draw[-,line width=1pt] (-0.3,-.3) to (.3,.4);
	\draw[-,line width=4pt,white] (0.3,-.3) to (-.3,.4);
  \draw[-,line width=1pt] (0.3,-.3) to (-.3,.4);
        \node at (-0.3,-.4) {$\scriptstyle a$};
        \node at (0.3,-.4) {$\scriptstyle b$};
     \draw(0.15,0.23)   \bdot;
\end{tikzpicture}
,  
%
    \qquad
\begin{tikzpicture}[baseline=-1mm, scale=.8, color=\clr]
	\draw[-,line width=1pt] (0.3,-.3) to (-.3,.4);
	\draw[-,line width=4pt,white] (-0.3,-.3) to (.3,.4);
	\draw[-,line width=1pt] (-0.3,-.3) to (.3,.4);
        \node at (-0.3,-.4) {$\scriptstyle a$};
        \node at (0.3,-.4) {$\scriptstyle b$};
     \draw(-0.15,0.23)   \bdot;
\end{tikzpicture}
  =
\begin{tikzpicture}[baseline=-1mm, scale=.8, color=\clr]
 \draw[-,line width=1pt] (-0.3,-.3) to (.3,.4);
	\draw[-,line width=4pt,white] (0.3,-.3) to (-.3,.4);
  \draw[-,line width=1pt] (0.3,-.3) to (-.3,.4);
        \node at (-0.3,-.4) {$\scriptstyle a$};
        \node at (0.3,-.4) {$\scriptstyle b$};
     \draw(0.15,-0.12)   \bdot;
\end{tikzpicture}
,
\\
 \label{dotmovesplits+merge}
\begin{tikzpicture}[baseline = -.5mm,scale=.8,color=\clr]
\draw[-,line width=1.5pt] (0.08,-.5) to (0.08,0.04);
\draw[-,line width=1pt] (0.34,.5) to (0.08,0);
\draw[-,line width=1pt] (-0.2,.5) to (0.08,0);
\node at (-0.22,.6) {$\scriptstyle a$};
\node at (0.36,.65) {$\scriptstyle b$};
\draw (0.08,-.2) \bdot;
%\draw (0.45,-.2) node{$\scriptstyle \omega_{a+b}$};
\end{tikzpicture} 
& =
\begin{tikzpicture}[baseline = -.5mm,scale=.8,color=\clr]
\draw[-,line width=1.5pt] (0.08,-.5) to (0.08,0.04);
\draw[-,line width=1pt] (0.34,.5) to (0.08,0);
\draw[-,line width=1pt] (-0.2,.5) to (0.08,0);
\node at (-0.22,.6) {$\scriptstyle a$};
\node at (0.36,.65) {$\scriptstyle b$};
\draw (-.05,.24) \bdot;
%\draw (-0.3,.2) node{$\scriptstyle \omega_{a}$};
%\draw (0.48,.2) node{$\scriptstyle \omega_{b}$};
\draw (.22,.24) \bdot;
\end{tikzpicture},
   \quad 
\begin{tikzpicture}[baseline = -.5mm, scale=.8, color=\clr]
\draw[-,line width=1pt] (0.3,-.5) to (0.08,0.04);
\draw[-,line width=1pt] (-0.2,-.5) to (0.08,0.04);
\draw[-,line width=1.5pt] (0.08,.6) to (0.08,0);
\node at (-0.22,-.6) {$\scriptstyle a$};
\node at (0.35,-.6) {$\scriptstyle b$};
\draw (0.08,.2) \bdot;
\end{tikzpicture}
 =~
\begin{tikzpicture}[baseline = -.5mm,scale=.8, color=\clr]
\draw[-,line width=1pt] (0.3,-.5) to (0.08,0.04);
\draw[-,line width=1pt] (-0.2,-.5) to (0.08,0.04);
\draw[-,line width=1.5pt] (0.08,.6) to (0.08,0);
\node at (-0.22,-.6) {$\scriptstyle a$};
\node at (0.35,-.6) {$\scriptstyle b$};
\draw (-.08,-.3) \bdot; \draw (.22,-.3) \bdot;
%\draw (-0.41,-.3) node{$\scriptstyle \omega_{a}$};
%\draw (0.48,-.3) node{$\scriptstyle \omega_{b}$};
\end{tikzpicture} .
\end{align}
\end{definition}
The following relations hold in $\qAW$:
\begin{align}
  \label{intergralballon}
\begin{tikzpicture}[baseline = 1.5mm, scale=.5, color=\clr]
\draw[-, line width=1.2pt] (0.5,2) to (0.5,2.5);
\draw[-, line width=1.2pt] (0.5,0) to (0.5,-.4);
\draw[-,thin]  (0.5,2) to[out=left,in=up] (-.5,1)
to[out=down,in=left] (0.5,0);
\draw[-,thin]  (0.5,2) to[out=left,in=up] (0,1)
 to[out=down,in=left] (0.5,0);      
\draw[-,thin] (0.5,0)to[out=right,in=down] (1.5,1)
to[out=up,in=right] (0.5,2);
\draw[-,thin] (0.5,0)to[out=right,in=down] (1,1) to[out=up,in=right] (0.5,2);
\node at (0.5,.7){$\scriptstyle \cdots$};
\draw (-0.5,1) \bdot; 
%\node at (-.8,1) {$\scriptstyle \omega_1$}; 
\draw (0,1) \bdot; 
%\node at (0.3,1) {$\scriptstyle \omega_1$};
\node at (0.5,-.6) {$\scriptstyle a$};
\draw (1,1) \bdot;
\draw (1.5,1) \bdot; 
%\node at (1.8,1) {$\scriptstyle \omega_1$};
\node at (-.22,0) {$\scriptstyle 1$};
\node at (1.2,0) {$\scriptstyle 1$};
\node at (.3,0.3) {$\scriptstyle 1$};
\node at (.7,0.3) {$\scriptstyle 1$};
\end{tikzpicture}
   & =
   [a]_q!  
\begin{tikzpicture}[baseline = 1.5mm, scale=.5, color=\clr]
\draw[-,line width=1.2pt] (0,-0.1) to[out=up, in=down] (0,1.4);
\draw(0,0.6) \bdot; 
%\draw (0.3,0.6) node {$\scriptstyle \omega_a$};
\node at (0,-.3) {$\scriptstyle a$};
\end{tikzpicture} ,
\qquad
\begin{tikzpicture}[baseline = 1.5mm, scale=.5, color=\clr]
\draw[-, line width=1.2pt] (0.5,2) to (0.5,2.5);
\draw[-, line width=1.2pt] (0.5,0) to (0.5,-.4);
\draw[-,thin]  (0.5,2) to[out=left,in=up] (-.5,1)
to[out=down,in=left] (0.5,0);
\draw[-,thin]  (0.5,2) to[out=left,in=up] (0,1)
 to[out=down,in=left] (0.5,0);      
\draw[-,thin] (0.5,0)to[out=right,in=down] (1.5,1)
to[out=up,in=right] (0.5,2);
\draw[-,thin] (0.5,0)to[out=right,in=down] (1,1) to[out=up,in=right] (0.5,2);
\node at (0.5,.7){$\scriptstyle \cdots$};
\draw (-0.5,1) \zdot; 
%\node at (-.8,1) {$\scriptstyle \omega_1$}; 
\draw (0,1) \zdot; 
%\node at (0.3,1) {$\scriptstyle \omega_1$};
\node at (0.5,-.6) {$\scriptstyle a$};
\draw (1,1) \zdot;
\draw (1.5,1) \zdot; 
%\node at (1.8,1) {$\scriptstyle \omega_1$};
\node at (-.22,0) {$\scriptstyle 1$};
\node at (1.2,0) {$\scriptstyle 1$};
\node at (.3,0.3) {$\scriptstyle 1$};
\node at (.7,0.3) {$\scriptstyle 1$};
\end{tikzpicture}
    ~=~ [a]_q!  
\begin{tikzpicture}[baseline = 1.5mm, scale=.5, color=\clr]
\draw[-,line width=1.2pt] (0,-0.1) to[out=up, in=down] (0,1.4);
\draw(0,0.6) \zdot; 
%\draw (0.3,0.6) node {$\scriptstyle \omega_a$};
\node at (0,-.3) {$\scriptstyle a$};
\end{tikzpicture} .
\end{align}
The $a$-fold merges (and splits) used in \eqref{intergralballon} are defined by using \eqref{webassoc} via iterated merges (and splits), and then the relations \eqref{intergralballon} follow from \eqref{webassoc}, \eqref{splitbinomial} and \eqref{dotmovesplits+merge}. 

We define the following morphisms in $\qAW$
\begin{equation}
\label{eq:wkdota}
\omega_{a,r} :=
\begin{tikzpicture}[baseline = -1mm,scale=.7,color=\clr]
\draw[-,line width=1.2pt] (0.08,-.7) to (0.08,.5);
\node at (.08,-.9) {$\scriptstyle a$};
\draw(0.08,-0.1) \bdot;
\draw(.5,0)node {$\scriptstyle \omega_r$};
\end{tikzpicture}
:=
\begin{tikzpicture}[baseline = -1mm,scale=.7,color=\clr]
\draw[-,line width=1.5pt] (0.08,-.8) to (0.08,-.5);
\draw[-,line width=1.5pt] (0.08,.3) to (0.08,.6);
\draw[-,line width=1pt] (0.1,-.51) to [out=45,in=-45] (0.1,.31);
\draw[-,line width=1pt] (0.06,-.51) to [out=135,in=-135] (0.06,.31);
\draw(-.1,-0.1) \bdot;
\node at (-.25,-.35) {$\scriptstyle r$};
\node at (.08,-1.1) {$\scriptstyle a$};
\end{tikzpicture} \; ,
\qquad
\omega_{a,-r} :=
\begin{tikzpicture}[baseline = -1mm,scale=.7,color=\clr]
\draw[-,line width=1.2pt] (0.08,-.7) to (0.08,.5);
\node at (.08,-.9) {$\scriptstyle a$};
\draw(0.08,-0.1) \bdot;
\draw(.6,0)node {$\scriptstyle \omega_{-r}$};
\end{tikzpicture}
:=
\begin{tikzpicture}[baseline = -1mm,scale=.7,color=\clr]
\draw[-,line width=1.2pt] (0.08,-.7) to (0.08,.5);
\node at (.08,-.9) {$\scriptstyle a$};
\draw(0.08,-0.1) \zdot;
\draw(.5,0)node {$\scriptstyle \omega_r$};
\end{tikzpicture}
:=
\begin{tikzpicture}[baseline = -1mm,scale=.7,color=\clr]
\draw[-,line width=1.5pt] (0.08,-.8) to (0.08,-.5);
\draw[-,line width=1.5pt] (0.08,.3) to (0.08,.6);
\draw[-,line width=1pt] (0.1,-.51) to [out=45,in=-45] (0.1,.31);
\draw[-,line width=1pt] (0.06,-.51) to [out=135,in=-135] (0.06,.31);
\draw(.27,-0.1) \zdot;
\node at (.4,-.35) {$\scriptstyle r$};
\node at (.08,-.9) {$\scriptstyle a$};
\end{tikzpicture}  
\qquad\quad (0 \le r\le a).
\end{equation}
In particular, we have 
$
\begin{tikzpicture}[baseline = 3pt, scale=0.4, color=\clr]
\draw[-,line width=1.2pt] (0,0) to[out=up, in=down] (0,1.4);
\draw(0,0.6) \bdot;
\draw(0.7,0.6)node {$\scriptstyle \omega_a$};
\node at (0,-.22) {$\scriptstyle a$};
\end{tikzpicture}
=\xdota$
and 
$\begin{tikzpicture}[baseline = 3pt, scale=0.4, color=\clr]
\draw[-,line width=1.2pt] (0,0) to[out=up, in=down] (0,1.4);
\draw(0,0.6) \zdot;
\draw(0.7,0.6)node {$\scriptstyle \omega_a$};
\node at (0,-.22) {$\scriptstyle a$};
\end{tikzpicture}
=\zdota$. 

% The following relations in $\qAW$ follow by \eqref{dotmovecrossing} and the inverse relations between solid and hollow dots:
% \begin{align}
% \label{eq:zdotmovecrossing}
%     \begin{tikzpicture}[baseline=-1mm, scale=.8, color=\clr]
%  \draw[-,line width=1pt] (-0.3,-.3) to (.3,.4);
% 	\draw[-,line width=4pt,white] (0.3,-.3) to (-.3,.4);
%   \draw[-,line width=1pt] (0.3,-.3) to (-.3,.4);
%         \node at (-0.3,-.4) {$\scriptstyle a$};
%         \node at (0.3,-.4) {$\scriptstyle b$};
%      \draw(-0.18,-0.15)   \zdot;
% \end{tikzpicture}
% =
% \begin{tikzpicture}[baseline=-1mm, scale=.8, color=\clr]
% 	\draw[-,line width=1pt] (0.3,-.3) to (-.3,.4);
% 	\draw[-,line width=4pt,white] (-0.3,-.3) to (.3,.4);
% 	\draw[-,line width=1pt] (-0.3,-.3) to (.3,.4);
%         \node at (-0.3,-.4) {$\scriptstyle a$};
%         \node at (0.3,-.4) {$\scriptstyle b$};
%      \draw(0.15,0.23)   \zdot;
% \end{tikzpicture} ,
%     \qquad
% \begin{tikzpicture}[baseline=-1mm, scale=.8, color=\clr]
%  \draw[-,line width=1pt] (-0.3,-.3) to (.3,.4);
% 	\draw[-,line width=4pt,white] (0.3,-.3) to (-.3,.4);
%   \draw[-,line width=1pt] (0.3,-.3) to (-.3,.4);
%         \node at (-0.3,-.4) {$\scriptstyle a$};
%         \node at (0.3,-.4) {$\scriptstyle b$};
%       \draw(-0.15,0.23)   \zdot;
% \end{tikzpicture} 
% =
% \begin{tikzpicture}[baseline=-1mm, scale=.8, color=\clr]
% 	\draw[-,line width=1pt] (0.3,-.3) to (-.3,.4);
% 	\draw[-,line width=4pt,white] (-0.3,-.3) to (.3,.4);
% 	\draw[-,line width=1pt] (-0.3,-.3) to (.3,.4);
%         \node at (-0.3,-.4) {$\scriptstyle a$};
%         \node at (0.3,-.4) {$\scriptstyle b$};
%     \draw(0.18,-0.15)   \zdot;
% \end{tikzpicture} .
% \end{align}

% Similarly, the following relations in $\qAW$, for $a,b\ge 1$, follow by \eqref{dotmovesplits+merge}:
% \begin{equation}
% \label{eq:dotSM2}
% \begin{tikzpicture}[baseline = -.5mm,scale=.8,color=\clr]
% \draw[-,line width=1.5pt] (0.08,-.5) to (0.08,0.04);
% \draw[-,line width=1pt] (0.34,.5) to (0.08,0);
% \draw[-,line width=1pt] (-0.2,.5) to (0.08,0);
% \node at (-0.22,.6) {$\scriptstyle a$};
% \node at (0.36,.65) {$\scriptstyle b$};
% \draw (0.08,-.2) \zdot;
% \end{tikzpicture} 
% =
% \begin{tikzpicture}[baseline = -.5mm,scale=.8,color=\clr]
% \draw[-,line width=1.5pt] (0.08,-.5) to (0.08,0.04);
% \draw[-,line width=1pt] (0.34,.5) to (0.08,0);
% \draw[-,line width=1pt] (-0.2,.5) to (0.08,0);
% \node at (-0.22,.6) {$\scriptstyle a$};
% \node at (0.36,.65) {$\scriptstyle b$};
% \draw (-.05,.24) \zdot;
% \draw (.22,.24) \zdot;
% \end{tikzpicture},
%    \qquad 
% \begin{tikzpicture}[baseline = -.5mm, scale=.8, color=\clr]
% \draw[-,line width=1pt] (0.3,-.5) to (0.08,0.04);
% \draw[-,line width=1pt] (-0.2,-.5) to (0.08,0.04);
% \draw[-,line width=1.5pt] (0.08,.6) to (0.08,0);
% \node at (-0.22,-.6) {$\scriptstyle a$};
% \node at (0.35,-.6) {$\scriptstyle b$};
% \draw (0.08,.2) \zdot;
% \end{tikzpicture}
%  =~
% \begin{tikzpicture}[baseline = -.5mm,scale=.8, color=\clr]
% \draw[-,line width=1pt] (0.3,-.5) to (0.08,0.04);
% \draw[-,line width=1pt] (-0.2,-.5) to (0.08,0.04);
% \draw[-,line width=1.5pt] (0.08,.6) to (0.08,0);
% \node at (-0.22,-.6) {$\scriptstyle a$};
% \node at (0.35,-.6) {$\scriptstyle b$};
% \draw (-.08,-.3) \zdot; \draw (.22,-.3) \zdot;
% \end{tikzpicture} .
% \end{equation}

%\subsection{The $q$-web category}

% To be consistent with our general terminology, the $q$-Schur category $q$-{\bf Schur} in \cite{Bru25} will be referred to as the {\em $q$-web category} and denoted by $\qW$. More explicitly, $\qW$ is the strict $\kk$-linear monoidal category generated by objects $a\in \mathbb Z_{\ge1}$. The object $a$ and its identity morphism  will be drawn as a vertical strand labeled by $a$:
% $
% \stra .  
% $
% The generating morphisms are the merges, splits and (thick) positive crossings depicted as \eqref{merge+split+crossing}, subject to the relations \eqref{webassoc}--\eqref{splitbinomial}, where $\crossingpos$ is invertible with inverse denoted by $\crossingneg$. 

% The following additional relations hold for $\qW$ (and hence also for $\qAW$) as they can be derived from the defining relations of $\qW$, for $a,b,c,d \geq 0$ 
% with $d \leq a$ and $c \leq b+d$: \red{do we really need them or all of them?}
% \begin{align}
% %\text{(square-switch relations)}
% %\notag \\
% %square switch
% \label{squareswitch}
% \begin{tikzpicture}[anchorbase,scale=.9, color=\clr]
% 	\draw[-,thick] (0,0) to (0,1);
% 	\draw[-,thick] (.015,0) to (0.015,.2) to (.57,.4) to (.57,.6)
%         to (.015,.8) to (.015,1);
% 	\draw[-,line width=1.2pt] (0.6,0) to (0.6,1);
%         \node at (0.6,-.2) {$\scriptstyle b$};
%         \node at (0,-.2) {$\scriptstyle a$};
%         \node at (0.3,.95) {$\scriptstyle c$};
%         \node at (0.3,.12) {$\scriptstyle d$};
% \end{tikzpicture}
% &=
% \!\!\!\sum_{s=\max(0,c-b)}^{\min(c,d)}
% \qbinomsm{a -b+c-\!d}{s} \ 
% \begin{tikzpicture}[anchorbase,scale=.9, color=\clr]
% 	\draw[-,line width=1.2pt] (0,0) to (0,1);
% 	\draw[-,thick] (0.8,0) to (0.8,.2) to (.03,.4) to (.03,.6)
%         to (.8,.8) to (.8,1);
% 	\draw[-,thin] (0.82,0) to (0.82,1);
%         \node at (0.81,-.15) {$\scriptstyle b$};
%         \node at (0,-.15) {$\scriptstyle a$};
%         \node at (0.4,.9) {$\scriptstyle d-s$};
%         \node at (0.4,.13) {$\scriptstyle c-s$};
% \end{tikzpicture},
% \\
% %squareswitch2
% \label{squareswitch2}
% \begin{tikzpicture}[anchorbase,scale=.9, color=\clr]
% 	\draw[-,thick] (0,0) to (0,1);
% 	\draw[-,thick] (-.015,0) to (-0.015,.2) to (-.57,.4) to (-.57,.6)
%         to (-.015,.8) to (-.015,1);
% 	\draw[-,line width=1.2pt] (-0.6,0) to (-0.6,1);
%         \node at (-0.6,-.15) {$\scriptstyle b$};
%         \node at (0,-.15) {$\scriptstyle a$};
%         \node at (-0.3,.88) {$\scriptstyle c$};
%         \node at (-0.3,.15) {$\scriptstyle d$};
% \end{tikzpicture}
% &=\!\!\!
% \sum_{s=\max(0,c-b)}^{\min(c,d)}
% \qbinomsm{a-b+c-d}{s} \ 
% \begin{tikzpicture}[anchorbase,scale=.9, color=\clr]
% 	\draw[-,line width=1.2pt] (0,0) to (0,1);
% 	\draw[-,thick] (-0.8,0) to (-0.8,.2) to (-.03,.4) to (-.03,.6)
%         to (-.8,.8) to (-.8,1);
% 	\draw[-,thin] (-0.82,0) to (-0.82,1);
%         \node at (-0.81,-.15) {$\scriptstyle b$};
%         \node at (0,-.15) {$\scriptstyle a$};
%         \node at (-0.4,.9) {$\scriptstyle d-s$};
%         \node at (-0.4,.13) {$\scriptstyle c-s$};
% \end{tikzpicture} ,
% \\ 
% %1cross+
% \label{cross via square}
% \begin{tikzpicture}[baseline=-1mm, scale=.9, color=\clr]
% 	\draw[-,line width=1pt] (0.3,-.3) to (-.3,.4);
% 	\draw[-,line width=4pt,white] (-0.3,-.3) to (.3,.4);
% 	\draw[-,line width=1pt] (-0.3,-.3) to (.3,.4);
%         \node at (-0.3,-.45) {$\scriptstyle a$};
%         \node at (0.3,-.45) {$\scriptstyle b$};
% \end{tikzpicture}
% & =\sum_{s=0}^{\min(a,b)}
% (-q)^{s}
% \begin{tikzpicture}[anchorbase,scale=.9, color=\clr]
% 	\draw[-,line width=1.2pt] (0,0) to (0,1);
% 	\draw[-,thick] (-0.8,0) to (-0.8,.2) to (-.03,.4) to (-.03,.6)
%         to (-.8,.8) to (-.8,1);
% 	\draw[-,thin] (-0.82,0) to (-0.82,1);
%         \node at (-0.81,-.15) {$\scriptstyle a$};
%         \node at (0,-.15) {$\scriptstyle b$};
%         \node at (-0.4,.9) {$\scriptstyle b-s$};
%         \node at (-0.4,.13) {$\scriptstyle a-s$};
% \end{tikzpicture}
% =
% \sum_{s=0}^{\min(a,b)}
% (-q)^{s}
% \begin{tikzpicture}[anchorbase,scale=.9, color=\clr]
% 	\draw[-,line width=1.2pt] (0,0) to (0,1);
% 	\draw[-,thick] (0.8,0) to (0.8,.2) to (.03,.4) to (.03,.6)
%         to (.8,.8) to (.8,1);
% 	\draw[-,thin] (0.82,0) to (0.82,1);
%         \node at (0.81,-.15) {$\scriptstyle b$};
%         \node at (0,-.15) {$\scriptstyle a$};
%         \node at (0.4,.9) {$\scriptstyle a-s$};
%         \node at (0.4,.13) {$\scriptstyle b-s$};
% \end{tikzpicture},
% \\   
% %2cross-
% \begin{tikzpicture}[baseline=-1mm, scale=.9, color=\clr]
%  \draw[-,line width=1pt] (-0.3,-.3) to (.3,.4);
% 	\draw[-,line width=4pt,white] (0.3,-.3) to (-.3,.4);
%  \draw[-,line width=1pt] (0.3,-.3) to (-.3,.4);
%         \node at (-0.3,-.45) {$\scriptstyle a$};
%         \node at (0.3,-.45) {$\scriptstyle b$};
% \end{tikzpicture}
% &=\sum_{s=0}^{\min(a,b)}
% (-q)^{-s}
% \begin{tikzpicture}[anchorbase,scale=.9, color=\clr]
% 	\draw[-,line width=1.2pt] (0,0) to (0,1);
% 	\draw[-,thick] (-0.8,0) to (-0.8,.2) to (-.03,.4) to (-.03,.6)
%         to (-.8,.8) to (-.8,1);
% 	\draw[-,thin] (-0.82,0) to (-0.82,1);
%         \node at (-0.81,-.15) {$\scriptstyle a$};
%         \node at (0,-.15) {$\scriptstyle b$};
%         \node at (-0.4,.9) {$\scriptstyle b-s$};
%         \node at (-0.4,.13) {$\scriptstyle a-s$};
% \end{tikzpicture}
% =
% \sum_{s=0}^{\min(a,b)}
% (-q)^{-s}
% \begin{tikzpicture}[anchorbase,scale=.9, color=\clr]
% 	\draw[-,line width=1.2pt] (0,0) to (0,1);
% 	\draw[-,thick] (0.8,0) to (0.8,.2) to (.03,.4) to (.03,.6)
%         to (.8,.8) to (.8,1);
% 	\draw[-,thin] (0.82,0) to (0.82,1);
%         \node at (0.81,-.15) {$\scriptstyle b$};
%         \node at (0,-.15) {$\scriptstyle a$};
%         \node at (0.4,.9) {$\scriptstyle a-s$};
%         \node at (0.4,.13) {$\scriptstyle b-s$};
% \end{tikzpicture} ,
% \end{align}


% \vspace{-4mm}

% \begin{align}
% %sliders
% \begin{tikzpicture}[anchorbase,scale=0.7, color=\clr]
% 	\draw[-,line width=1pt] (0.4,0) to (-0.6,1);
% 	\draw[-,line width=4pt,white] (0.1,0) to (0.1,.6) to (.5,1);
% 	\draw[-,line width=1pt] (0.1,0) to (0.1,.6) to (.5,1);
% 	\draw[-,line width=1pt] (0.08,0) to (0.08,1);
%         \node at (0.6,1.13) {$\scriptstyle c$};
%         \node at (0.1,1.19) {$\scriptstyle b$};
%         \node at (-0.65,1.14) {$\scriptstyle a$};
% \end{tikzpicture}
% &=
% \begin{tikzpicture}[anchorbase,scale=0.7, color=\clr]
% 	\draw[-,line width=1pt] (0.7,0) to (-0.3,1);
% 	\draw[-,line width=4pt,white] (0.1,0) to (0.1,.2) to (.9,1);
% 	\draw[-,line width=3pt,white] (0.08,0) to (0.08,1);
% 	\draw[-,line width=1pt] (0.1,0) to (0.1,.2) to (.9,1);
% 	\draw[-,line width=1pt] (0.08,0) to (0.08,1);
%         \node at (0.9,1.13) {$\scriptstyle c$};
%         \node at (0.1,1.19) {$\scriptstyle b$};
%         \node at (-0.4,1.13) {$\scriptstyle a$};
% \end{tikzpicture},&
% \begin{tikzpicture}[anchorbase,scale=0.7, color=\clr]
% 	\draw[-,line width=1pt] (-0.08,0) to (-0.08,1);
% 	\draw[-,line width=1pt] (-0.1,0) to (-0.1,.6) to (-.5,1);
% 	\draw[-,line width=4pt,white] (-0.4,0) to (0.6,1);
% 	\draw[-,line width=1pt] (-0.4,0) to (0.6,1);
%         \node at (0.7,1.13) {$\scriptstyle c$};
%         \node at (-0.1,1.19) {$\scriptstyle b$};
%         \node at (-0.6,1.13) {$\scriptstyle a$};
% \end{tikzpicture}
% &=
% \begin{tikzpicture}[anchorbase,scale=0.7,  color=\clr]
% 	\draw[-,line width=1pt] (-0.08,0) to (-0.08,1);
% 	\draw[-,line width=1pt] (-0.1,0) to (-0.1,.2) to (-.9,1);
% 	\draw[-,line width=4pt,white] (-0.7,0) to (0.3,1);
% 	\draw[-,line width=1pt] (-0.7,0) to (0.3,1);
%         \node at (0.4,1.13) {$\scriptstyle c$};
%         \node at (-0.1,1.19) {$\scriptstyle b$};
%         \node at (-0.95,1.13) {$\scriptstyle a$};
% \end{tikzpicture},&
% \:\begin{tikzpicture}[baseline=-3.3mm,scale=0.7, color=\clr]
% 	\draw[-,line width=1pt] (0.08,0) to (0.08,-1);
% 	\draw[-,line width=1pt] (0.1,0) to (0.1,-.6) to (.5,-1);
% 	\draw[-,line width=4pt,white] (0.4,0) to (-0.6,-1);
% 	\draw[-,line width=1pt] (0.4,0) to (-0.6,-1);
%         \node at (0.6,-1.13) {$\scriptstyle c$};
%         \node at (0.05,-1.16) {$\scriptstyle b$};
%         \node at (-0.6,-1.13) {$\scriptstyle a$};
% \end{tikzpicture}
% &=
% \begin{tikzpicture}[baseline=-3.3mm,scale=0.7, color=\clr]
% 	\draw[-,line width=1pt] (0.08,0) to (0.08,-1);
% 	\draw[-,line width=1pt] (0.1,0) to (0.1,-.2) to (.9,-1);
% 	\draw[-,line width=4pt, white] (0.7,0) to (-0.3,-1);
% 	\draw[-,line width=1pt] (0.7,0) to (-0.3,-1);
%         \node at (1,-1.13) {$\scriptstyle c$};
%         \node at (0.19,-1.16) {$\scriptstyle b$};
%         \node at (-0.4,-1.13) {$\scriptstyle a$};
% \end{tikzpicture},\:&
% \begin{tikzpicture}[baseline=-3.3mm,scale=0.7, color=\clr]
% 	\draw[-,line width=1pt] (-0.4,0) to (0.6,-1);
% 	\draw[-,line width=4pt,white] (-0.1,0) to (-0.1,-.6) to (-.5,-1);
% 	\draw[-,line width=1pt] (-0.1,0) to (-0.1,-.6) to (-.5,-1);
% 	\draw[-,line width=1pt] (-0.08,0) to (-0.08,-1);
%         \node at (0.6,-1.13) {$\scriptstyle c$};
%         \node at (0.01,-1.16) {$\scriptstyle b$};
%         \node at (-0.6,-1.13) {$\scriptstyle a$};
% \end{tikzpicture}
% &=
% \begin{tikzpicture}[baseline=-3.3mm,scale=0.7, color=\clr]
% 	\draw[-,line width=1pt] (-0.7,0) to (0.3,-1);
% 	\draw[-,line width=4pt,white] (-0.1,0) to (-0.1,-.2) to (-.9,-1);
% 	\draw[-,line width=3pt,white] (-0.08,0) to (-0.08,-1);
% 	\draw[-,line width=1pt] (-0.1,0) to (-0.1,-.2) to (-.9,-1);
% 	\draw[-,line width=1pt] (-0.08,0) to (-0.08,-1);
%         \node at (0.34,-1.13) {$\scriptstyle c$};
%         \node at (-0.12,-1.16) {$\scriptstyle b$};
%         \node at (-0.95,-1.13) {$\scriptstyle a$};
% \end{tikzpicture},
% \label{sliders}
% \end{align}

% \vspace{-4mm}

% \begin{align}
% % mergecross, crosssplit, 2cross+-
% \begin{tikzpicture}[anchorbase,scale=.6, color=\clr]
% \draw[-,line width=1pt] (-.2,-.8) to [out=45,in=-45] (0.1,.31);
% \draw[-,line width=4.3pt, white] (.36,-.8) to [out=135,in=-135] (0.06,.31);
% \draw[-,line width=1pt] (.36,-.8) to [out=135,in=-135] (0.06,.31);
% 	\draw[-,line width=1.5pt] (0.08,.3) to (0.08,.5);
%         \node at (-.3,-1) {$\scriptstyle a$};
%         \node at (.45,-1) {$\scriptstyle b$};
% \end{tikzpicture}
% &=
% q^{ab}\!\begin{tikzpicture}[anchorbase,scale=.6, color=\clr]
% 	\draw[-,line width=1.5pt] (0.08,.1) to (0.08,.5);
% \draw[-,line width=1pt] (.46,-.8) to [out=100,in=-45] (0.1,.11);
% \draw[-,line width=1pt] (-.3,-.8) to [out=80,in=-135] (0.06,.11);
%         \node at (-.3,-1) {$\scriptstyle a$};
%         \node at (.43,-1) {$\scriptstyle b$};
% \end{tikzpicture},
% &
% \begin{tikzpicture}[baseline=.3mm,scale=.6, color=\clr]
% \draw[-,line width=1pt] (.36,.8) to [out=-135,in=135] (0.06,-.31);
% \draw[-,line width=4.3pt,white] (-.2,.8) to [out=-45,in=45] (0.1,-.31);
% \draw[-,line width=1pt] (-.2,.8) to [out=-45,in=45] (0.1,-.31);
% 	\draw[-,line width=1.5pt] (0.08,-.3) to (0.08,-.5);
%         \node at (-.3,1) {$\scriptstyle a$};
%         \node at (.45,1) {$\scriptstyle b$};
% \end{tikzpicture}
% &=q^{ab}
% \begin{tikzpicture}[baseline=.3mm,scale=.6, color=\clr]
% 	\draw[-,line width=1.5pt] (0.08,-.1) to (0.08,-.5);
% \draw[-,line width=1pt] (.46,.8) to [out=-100,in=45] (0.1,-.11);
% \draw[-,line width=1pt] (-.3,.8) to [out=-80,in=135] (0.06,-.11);
%         \node at (-.3,1) {$\scriptstyle a$};
%         \node at (.43,1) {$\scriptstyle b$};
% \end{tikzpicture},
% \label{unfold}
% &
% \begin{tikzpicture}[baseline = -1mm,scale=0.7, color=\clr]
% 	\draw[-,line width=1pt] (-0.28,0) to[out=90,in=-90] (0.28,.6);
% 	\draw[-,line width=1pt] (0.28,-.6) to[out=90,in=-90] (-0.28,0);
% 	\draw[-,line width=4pt,white] (0.28,0) to[out=90,in=-90] (-0.28,.6);
% 	\draw[-,line width=4pt,white] (-0.28,-.6) to[out=90,in=-90] (0.28,0);
% 	\draw[-,line width=1pt] (0.28,0) to[out=90,in=-90] (-0.28,.6);
% 	\draw[-,line width=1pt] (-0.28,-.6) to[out=90,in=-90] (0.28,0);
%         \node at (0.3,-.8) {$\scriptstyle b$};
%         \node at (-0.3,-.8) {$\scriptstyle a$};
% \end{tikzpicture}
% &=
% \begin{tikzpicture}[baseline = -1mm,scale=0.7, color=\clr]
% 	\draw[-,line width=1pt] (0.3,-.6) to (0.3,.6);
% 	\draw[-,line width=1pt] (-0.2,-.6) to (-0.2,.6);
%         \node at (0.3,-.8) {$\scriptstyle b$};
%         \node at (-0.2,-.8) {$\scriptstyle a$};
% \end{tikzpicture}\ ,&
% \begin{tikzpicture}[baseline = -1mm,scale=0.7, color=\clr]
% 	\draw[-,line width=1pt] (0.28,0) to[out=90,in=-90] (-0.28,.6);
% 	\draw[-,line width=1pt] (-0.28,-.6) to[out=90,in=-90] (0.28,0);
% 	\draw[-,line width=4pt,white] (-0.28,0) to[out=90,in=-90] (0.28,.6);
% 	\draw[-,line width=4pt,white] (0.28,-.6) to[out=90,in=-90] (-0.28,0);
% 	\draw[-,line width=1pt] (-0.28,0) to[out=90,in=-90] (0.28,.6);
% 	\draw[-,line width=1pt] (0.28,-.6) to[out=90,in=-90] (-0.28,0);
%         \node at (0.3,-.8) {$\scriptstyle a$};
%         \node at (-0.3,-.8) {$\scriptstyle b$};
% \end{tikzpicture}
% &=
% \begin{tikzpicture}[baseline = -1mm,scale=0.7, color=\clr]
% 	\draw[-,line width=1pt] (0.2,-.6) to (0.2,.6);
% 	\draw[-,line width=1pt] (-0.3,-.6) to (-0.3,.6);
%         \node at (0.2,-.8) {$\scriptstyle a$};
%         \node at (-0.3,-.8) {$\scriptstyle b$};
% \end{tikzpicture}\:,
% \end{align}

% \vspace{-4mm}

% \begin{align}
% %2cross++
% \begin{tikzpicture}[anchorbase,scale=0.7, color=\clr]
% 	\draw[-,line width=1pt] (0.28,0) to[out=90,in=-90] (-0.28,.6);
% 	\draw[-,line width=4pt,white] (-0.28,0) to[out=90,in=-90] (0.28,.6);
% 	\draw[-,line width=1pt] (0.28,-.6) to[out=90,in=-90] (-0.28,0);
% 	\draw[-,line width=4pt,white] (-0.28,-.6) to[out=90,in=-90] (0.28,0);
% 	\draw[-,line width=1pt] (-0.28,-.6) to[out=90,in=-90] (0.28,0);
% 	\draw[-,line width=1pt] (-0.28,0) to[out=90,in=-90] (0.28,.6);
%         \node at (0.3,-.8) {$\scriptstyle b$};
%         \node at (-0.3,-.8) {$\scriptstyle a$};
% \end{tikzpicture}
% =
% \sum_{s=0}^{\min(a,b)}
% q^{-s(s-1)/2 -s(a+b-2s)}(q^{-1}-q)^s [s]!
% \begin{tikzpicture}[anchorbase,scale=.8, color=\clr]
% 	\draw[-,thick] (0.58,0) to (0.58,.2) to (.02,.8) to (.02,1);
% 	\draw[-,line width=4pt,white] (0.02,0) to (0.02,.2) to (.58,.8) to (.58,1);
% 	\draw[-,thick] (0.02,0) to (0.02,.2) to (.58,.8) to (.58,1);
% 	\draw[-,thin] (0,0) to (0,1);
% 	\draw[-,thin] (0.6,0) to (0.6,1);
%         \node at (0,-.18) {$\scriptstyle a$};
%         \node at (0.6,-.18) {$\scriptstyle b$};
%         \node at (-0.4,.5) {$\scriptstyle a-s$};
%         \node at (1.05,.5) {$\scriptstyle b-s$};
% \end{tikzpicture} ,
% \label{cross2}
% \end{align}
% % 
% \vspace{-2mm}
% \begin{align}
% %braid
% \begin{tikzpicture}[baseline = -1mm,scale=0.7, color=\clr]
% 	\draw[-,line width=1pt] (0.45,-.6) to (-0.45,.6);
%         \draw[-,line width=1pt] (0,-.6) to[out=90,in=-90] (-.45,0);
%         \draw[-,line width=4pt,white] (-0.45,0) to[out=90,in=-90] (0,0.6);
%         \draw[-,line width=1pt] (-0.45,0) to[out=90,in=-90] (0,0.6);
% 	\draw[-,line width=4pt,white] (0.45,.6) to (-0.45,-.6);
% 	\draw[-,line width=1pt] (0.45,.6) to (-0.45,-.6);
%         \node at (0,-.85) {$\scriptstyle b$};
%         \node at (0.5,-.85) {$\scriptstyle c$};
%         \node at (-0.5,-.85) {$\scriptstyle a$};
% \end{tikzpicture}
% &=
% \begin{tikzpicture}[baseline = -1mm,scale=0.7, color=\clr]
% 	\draw[-,line width=1pt] (0.45,-.6) to (-0.45,.6);
%         \draw[-,line width=1pt] (0.45,0) to[out=90,in=-90] (0,0.6);
%         \draw[-,line width=4pt,white] (0,-.6) to[out=90,in=-90] (.45,0);
%         \draw[-,line width=1pt] (0,-.6) to[out=90,in=-90] (.45,0);
% 	\draw[-,line width=4pt,white] (0.45,.6) to (-0.45,-.6);
% 	\draw[-,line width=1pt] (0.45,.6) to (-0.45,-.6);
%         \node at (0,-.85) {$\scriptstyle b$};
%         \node at (0.5,-.85) {$\scriptstyle c$};
%         \node at (-0.5,-.85) {$\scriptstyle a$};
% \end{tikzpicture} .
% \label{braid}
% \end{align}
% For details also see \cite{CKM},\cite[(6.11)--(6.17), Cor. 6.2]{Bru25}.
% Note that \eqref{cross2} in case $a=b=1$ is simply the standard quadratic relation for Hecke algebras:
% \begin{align} \label{Hecke}
%    \Big( 
%       \begin{tikzpicture}[baseline=-1mm, scale=.8, color=\clr]
% 	\draw[-,thick] (0.3,-.3) to (-.3,.4);
% 	\draw[-,line width=4pt,white] (-0.3,-.3) to (.3,.4);
% 	\draw[-,thick] (-0.3,-.3) to (.3,.4);
%         \node at (-0.3,-.45) {$\scriptstyle 1$};
%         \node at (0.26,-.45) {$\scriptstyle 1$};
% \end{tikzpicture}
%    \Big)^2 = (q^{-1} -q) 
%    \begin{tikzpicture}[baseline=-1mm, scale=.8, color=\clr]
% 	\draw[-,thick] (0.3,-.3) to (-.3,.4);
% 	\draw[-,line width=4pt,white] (-0.3,-.3) to (.3,.4);
% 	\draw[-,thick] (-0.3,-.3) to (.3,.4);
%         \node at (-0.3,-.45) {$\scriptstyle 1$};
%         \node at (0.26,-.45) {$\scriptstyle 1$};
% \end{tikzpicture}
%    +
%    \begin{tikzpicture}[baseline = -1mm,scale=0.6, color=\clr]
% 	\draw[-,thick] (0.2,-.6) to (0.2,.6);
% 	\draw[-,thick] (-0.3,-.6) to (-0.3,.6);
%         \node at (0.2,-.85) {$\scriptstyle 1$};
%         \node at (-0.3,-.85) {$\scriptstyle 1$};
% \end{tikzpicture} . 
% \end{align}
% It is not always necessary to include $\crossingpos$ as generating morphisms in $\qW$ and $\qAW$ (as it can be expressed in terms of merges and splits as in \eqref{cross via square}) if one includes additional relations  \eqref{squareswitch}--\eqref{squareswitch2} as defining relations. However, including $\crossingpos$ helps to discover and simplify the presentation.



%\subsection{Path algebras of the affine web category and its subcategories}

For any small $\kk$-linear category $\mathcal C$,  its path algebra is the locally unital algebra  $\oplus_{i,j\in \text{ob}\mathcal C}\Hom_{\mathcal C}(i,j)$ with multiplication induced by composition of morphisms. 
%Sometimes, we also use $\mathcal{ C}$ itself to denote its path algebra if there is no confusion.

Recall a (resp. strict) composition $\lambda=(\lambda_1,\ldots,\lambda_k)$ of $n$ is a tuple of positive (resp. non-negative) integers such that $\sum_{k}\lambda_k=n$; denote $|\lambda|=n$. Let $\Lambda(n)$ (resp. $\Lambda_\st(n)$) be the set of all (resp. strict) compositions of $n$.
The set of objects of $\qAW$ can be identified with the set of  all strict compositions (including $\emptyset$ which is identified with the unit object).   Let $\qAW(n) $ be the full subcategory 
of $\qAW$ with object set $\Lambda_\st(n)$. Denote by  $\hwS(n)$ and $\hwS$ the path algebras of $\qAW(n) $ and $\qAW$, respectively. Then $\hwS=\oplus_{n\in \N}\hwS(n)$. 
Moreover, it is proved in \cite{SSW25} that $\hwS(n)$ is Morita equivalent to the usual affine $q$-Schur algebra $\ASch(n,n)$; cf. \eqref{def:affineSchur}.

Let $\qAH$ be the full ``thin strand" subcategory of $\qAW$ with object set $\{(1^n)\mid n\in \N\}$. Then by the basis theorem given in \cite{SSW25}, $\qAH$ is the well-known strict monoidal category for the affine Hecke algebras, i.e., the endomorphism algebra $\End_{\qAH} ((1^n)) =\AHe(n)$. Here the affine Hecke algebra $\AHe(n)$ of type $GL_n$  is the unital associative $\kk$-algebra generated by $H_i^{\pm 1}$, $X_j^{\pm 1}$, for $1\le i\le n-1$ and $1\le j\le n$, subject to the relations:
\begin{align*}
    H_iH_i^{-1} &=1 =H_i^{-1}H_i, \qquad  (H_i+q)(H_i-q^{-1}) =0,
    \\
    H_iH_{i+1}H_i &=H_{i+1}H_iH_{i+1}, \qquad H_iH_j =H_jH_i \quad (|i-j|>1),
    \\
     X_iX_i^{-1} &=1 =X_i^{-1}X_i, \qquad\;  X_iX_j =X_jX_i,
    \\
    H_iX_iH_i &=X_{i+1}, \qquad\qquad\quad H_iX_j =X_jH_i  \quad (j\neq i,i+1).
\end{align*}
Diagrammatically, $H_i$ is the positive crossing between the $i$th and $(i+1)$st strands and $X_i$ is the thin dot on the $i$th strand. Moreover, by the basis theorem of the affine  $q$-webs in \cite{SSW25}, 
\begin{equation}
\label{equ:affwebhecke}
    \AHe(n)\cong 1_{(1^n)} \hwS(n) 1_{(1^n)}.
\end{equation} 


\subsection{The center of $\qAW$}

The center $Z\big(\qAW(n)\big)$ of the small linear category $\qAW(n) $ is the commutative subalgebra of $\Pi_{\mu \in \Lambda_\st(n)} 1_\mu \qAW(n) 1_\mu$
consists of $(z_\mu)_{\mu\in \Lambda_\st(n)}$ such that $\theta z_\lambda =z_\mu \theta
$ for any $ \theta\in 1_\mu \qAW(n)1_\lambda$.
Then the center $\cZ_{\hwS(n)}$ of the algebra $\hwS(n)$ is 
  \[
  \cZ_{\hwS(n)}= \Big\{\sum_{\mu\in \Lambda_\st(n)}z_\mu\mid (z_\mu)_{\mu\in\Lambda_\st(n)}\in \cZ\big(\qAW(n)\big) \Big\}.
  \]
  The following is well known and it follows from the double centralizer property for $\hwS(n)$ and $\AHe(n)$; cf. \eqref{T:module}--\eqref{def:affineSchur}, for $N\ge n$.
\begin{proposition}
\label{prop:center}
    We have $\cZ_{\hwS(n)}\cong \cZ_{\AHe(n)}= \kk[X_1^{\pm},\ldots,X_n^{\pm}]^{\mathfrak S_n}$, where the isomorphism is given by $z=\sum_\mu z_\mu \mapsto z_{(1^n)}$.
\end{proposition}

\begin{example}
    A central element $z \in \cZ_{\hwS(n)}$ is determined by its thin strand summand $z_{(1^n)}$. Take 
$z_{(1^n)}=q^{n-1}(X_1+X_2+\ldots+ X_n)$. Let     $z_{\lambda}=\sum_{i=1}^kL^\lambda_i$, where    
    $$L^\lambda_i=  q^{n-\lambda_i} \begin{tikzpicture}[baseline = 3pt, scale=0.5, color=\clr]
\draw[-,line width=1pt] (-1,-.2) to (-1,1.2);
\node at (-1,-.5) {$\scriptstyle \lambda_1$};
\node at (-.4,.4) {$\ldots$};
%\draw[-,line width=1pt] (.5,-.2) to (.5,1.2);
%\node at (.5,-.5) {$\scriptstyle 1$};
\draw[-,line width=1pt] (1,-.2) to (1,1.2);
\draw(1,0.5) \bdot; \node at (1.6,0.6) {$\scriptstyle \omega_{1}$};
\node at (1,-.5) {$\scriptstyle \lambda_{i}$};
%\draw[-,line width=1pt] (1.5,-.2) to (1.5,1.2);
%\node at (1.5,-.5) {$\scriptstyle 1$};
\node at (2.1,.3) {$\ldots$};
\draw[-,line width=1pt] (3,-.2) to (3,1.2);
\node at (3,-.5) {$\scriptstyle \lambda_k$};
\end{tikzpicture}  $$ 
for any $\lambda=(\lambda_1,\lambda_2,\ldots, \lambda_k)\in \Lambda_\st(n)$.
Then $z=\sum _{\lambda\in\Lambda_\st(n)}z_\lambda$.
\end{example}


\subsection{Restriction functors for $\qAW$}

For any $a\le n$, denote
\begin{align} \label{eq:1an}
\begin{split}
\Lambda^a(n) &=\{\mu\in \Lambda_\st(n)\mid \mu=(\mu_1,\mu_2,\ldots,\mu_k) \text{ such that } \mu_k=a\},
\\
\mathbf{1}_{a,n} &= \sum_{\mu\in \Lambda^a(n)}1_\mu. 
\end{split}
\end{align}
Then $\mathbf{1}_{a,n}$ is an idempotent of $\hwS(n)$. Moreover, by the basis theorem of $\qAW$ in \cite[Theorem 4.16]{SSW25},  there is an obvious embedding
\begin{equation}
\label{equ:embedding}
 \iota_{a,n}: \hwS(n-a) \longrightarrow \mathbf{1}_{a,n}\hwS(n)\mathbf{1}_{a,n}   
\end{equation}
sending any diagram $D$ to $D\stra$, the diagram obtained from $D$ by adding the  strand with thickness $a$ on the right.  
Then via the embedding $\iota_{a,n}$ we have a restriction functor
\[ \text{Res}^{n}_{n-a}=\mathbf{1}_{a,n}\hwS(n)\otimes _{\hwS(n)}?: \hwS(n)\modf\longrightarrow \hwS(n-a)\modf .\]

For any $f(X_1,\ldots,X_n)\in \kk[X_1^\pm,\ldots,X_n^\pm]^{\mathfrak S_n}$, let $z_f\in \cZ_{\hwS(n)}$ be the inverse of the isomorphism given in Proposition \ref{prop:center}. Given $\mathbf b=(b_1,b_2,\ldots,b_n)\in \kk^n$, we define the central character 
\[ 
\chi_{\mathbf b}: \cZ_{\hwS(n)}\longrightarrow \kk , \quad z_f\mapsto f(b_1,\ldots,b_n).
\]
Then $\chi_{\mathbf a}=\chi_{\mathbf b}$ if and only if $\mathbf a\sim \mathbf b$, i.e., $\mathbf a$ and $\mathbf b$ lie in the same orbit of $\mathfrak S_n$.

For any finite dimensional $M\in \hwS(n)\modf$ and $\gamma\in \kk^n/\sim$, let 
\[
M[\gamma]= \{v\in M\mid (z-\chi_{\gamma}(z))^kv=0 \text{ for } z\in \cZ_{\hwS(n)}, k\gg0\}. 
\]

For any $\gamma\in \kk^n/\sim$, let 
$\hwS(n)\modf_\gamma$ be the full category of $\hwS(n)\modf$ consisting of all modules $M$ with $M[\gamma]=M$.
Then $ \hwS(n)\modf=\oplus_{\gamma\in \kk^n/\sim}\hwS(n)\modf_\gamma $. 
 
For any multiset $\bfc=\{c_1,\ldots,c_a\}$ in $\kk$, we denote $|\bfc|=a$, and define the $\bfc$-restriction functor 
\begin{align}  \label{iRes_affine_n}
\bfc\text{-Res}^n_{n-|\bfc|}: \hwS(n)\modf\longrightarrow \hwS(n-a)\modf
\end{align}
which sends $M \in \hwS(n)\modf_\gamma$,
for each $\gamma$, to $ \text{ Proj}_{\gamma\setminus\bfc}\circ\text{Res}^n_{n-a}(M)$, where $ \text{ Proj}_{\gamma\setminus \bfc}$ denotes the projection from $\hwS(n-a)\modf $ to $\hwS(n-a)\modf_{\gamma\setminus \bfc}$, and  $\gamma\setminus \bfc $ is the orbit obtained from $\gamma$ by deleting $\bfc$ from $\gamma$.
 This leads to a functor 
\begin{align}  \label{iRes_affine}
\bfc\text{-Res} := \sum_{n} \bfc\text{-Res}^n_{n-|\bfc|}: \hwS\modf \longrightarrow \hwS\modf, 
\end{align}
which in turn gives rise to a linear operator $f_\bfc$ on $[\hwS\modf]$.


% \subsection{Category $\qAWC$ over $\C(q)$}

% We now introduce a $\C(q)$-linear monoidal category $\qAWC$, a simpler variant of $\qAW$. 

% \begin{definition}
%   \label{def:qAWebC}
% The affine $q$-web category $\qAWC$ is the strict $\C(q)$-linear monoidal category generated by objects $a\in \mathbb Z_{\ge1}$. The object $a$ and its identity morphism  will be drawn as a vertical strand labeled by $a$, i.e. 
% $
% \stra .  
% $
% The generating morphisms are the merges, splits and (thick) positive crossings depicted as
% \begin{align}
% \label{merge+split+crossingC}
% \merge 
% &:(a,b) \rightarrow (a+b),&
% \splits
% &:(a+b)\rightarrow (a,b),&
% \crossingpos
% &:(a,b) \rightarrow (b,a),
% \end{align}
% and 
% \begin{equation}
% \label{dotgenerator1}
% \dotgen, 
% \end{equation}
% for $a,b \in \Z_{\ge 1}$, where $\crossingpos$ is invertible with inverse being the negative crossing denoted $\crossingneg := \crossingpos^{-1}$ and $\dotgen$ is invertible with inverse $\zdotgen$, 
% subject to the following relations \eqref{webassocC}--\eqref{dotmovecrossingC}, for $a,b,c,d \in \Z_{\ge 1}$ with $d-a=c-b$:
% \begin{align}
% \label{webassocC}
% \begin{tikzpicture}[baseline = 0,color=\clr]
% 	\draw[-,thick] (0.35,-.3) to (0.08,0.14);
% 	\draw[-,thick] (0.1,-.3) to (-0.04,-0.06);
% 	\draw[-,line width=1pt] (0.085,.14) to (-0.035,-0.06);
% 	\draw[-,thick] (-0.2,-.3) to (0.07,0.14);
% 	\draw[-,line width=1.5pt] (0.08,.45) to (0.08,.1);
%         \node at (0.45,-.41) {$\scriptstyle c$};
%         \node at (0.07,-.4) {$\scriptstyle b$};
%         \node at (-0.28,-.41) {$\scriptstyle a$};
% \end{tikzpicture}
% &=
% \begin{tikzpicture}[baseline = 0, color=\clr]
% 	\draw[-,thick] (0.36,-.3) to (0.09,0.14);
% 	\draw[-,thick] (0.06,-.3) to (0.2,-.05);
% 	\draw[-,line width=1pt] (0.07,.14) to (0.19,-.06);
% 	\draw[-,thick] (-0.19,-.3) to (0.08,0.14);
% 	\draw[-,line width=1.5pt] (0.08,.45) to (0.08,.1);
%         \node at (0.45,-.41) {$\scriptstyle c$};
%         \node at (0.04,-.4) {$\scriptstyle b$};
%         \node at (-0.28,-.41) {$\scriptstyle a$};
% \end{tikzpicture}\:,
% \qquad
% \begin{tikzpicture}[baseline = -1mm, color=\clr]
% 	\draw[-,thick] (0.35,.3) to (0.08,-0.14);
% 	\draw[-,thick] (0.1,.3) to (-0.04,0.06);
% 	\draw[-,line width=1pt] (0.085,-.14) to (-0.035,0.06);
% 	\draw[-,thick] (-0.2,.3) to (0.07,-0.14);
% 	\draw[-,line width=1.5pt] (0.08,-.45) to (0.08,-.1);
%         \node at (0.45,.4) {$\scriptstyle c$};
%         \node at (0.07,.42) {$\scriptstyle b$};
%         \node at (-0.28,.4) {$\scriptstyle a$};
% \end{tikzpicture}
% =\begin{tikzpicture}[baseline = -1mm, color=\clr]
% 	\draw[-,thick] (0.36,.3) to (0.09,-0.14);
% 	\draw[-,thick] (0.06,.3) to (0.2,.05);
% 	\draw[-,line width=1pt] (0.07,-.14) to (0.19,.06);
% 	\draw[-,thick] (-0.19,.3) to (0.08,-0.14);
% 	\draw[-,line width=1.5pt] (0.08,-.45) to (0.08,-.1);
%         \node at (0.45,.4) {$\scriptstyle c$};
%         \node at (0.07,.42) {$\scriptstyle b$};
%         \node at (-0.28,.4) {$\scriptstyle a$};
% \end{tikzpicture}\:,
% \\
% \label{mergesplitC}
% \begin{tikzpicture}[baseline = 7.5pt,scale=.8, color=\clr]
% 	\draw[-,line width=1pt] (0,0) to (.28,.3) to (.28,.7) to (0,1);
% 	\draw[-,line width=1pt] (.6,0) to (.31,.3) to (.31,.7) to (.6,1);
%         \node at (0,1.13) {$\scriptstyle b$};
%         \node at (0.63,1.13) {$\scriptstyle d$};
%         \node at (0,-.1) {$\scriptstyle a$};
%         \node at (0.63,-.1) {$\scriptstyle c$};
% \end{tikzpicture}
% &=
% \sum_{\substack{0 \leq s \leq \min(a,b)\\0 \leq t \leq \min(c,d)\\t-s=d-a}}
% q^{st}  
% \begin{tikzpicture}[baseline = 7.5pt,scale=.8, color=\clr]
% \draw[-,thick] (0.58,0) to (0.58,.2) to (.02,.8) to (.02,1);
% 	\draw[-,line width=4pt,white] (0.02,0) to (0.02,.2) to (.58,.8) to (.58,1);
% 	\draw[-,thick] (0.02,0) to (0.02,.2) to (.58,.8) to (.58,1);
% 	\draw[-,thick] (0,0) to (0,1);
% 	\draw[-,thick] (0.6,0) to (0.6,1);
%         \node at (0,1.13) {$\scriptstyle c$};
%         \node at (0.6,1.13) {$\scriptstyle d$};
%         \node at (0,-.1) {$\scriptstyle a$};
%         \node at (0.57,-.1) {$\scriptstyle b$};
%         \node at (-0.1,.5) {$\scriptstyle s$};
%         \node at (0.75,.5) {$\scriptstyle t$};
% \end{tikzpicture},
% \\
%   \label{splitbinomialC}
% \begin{tikzpicture}[baseline = -1mm,scale=.8,color=\clr]
% \draw[-,line width=1.5pt] (0.08,-.8) to (0.08,-.5);
% \draw[-,line width=1.5pt] (0.08,.3) to (0.08,.6);
% \draw[-,thick] (0.1,-.51) to [out=45,in=-45] (0.1,.31);
% \draw[-,thick] (0.06,-.51) to [out=135,in=-135] (0.06,.31);
% \node at (-.33,-.05) {$\scriptstyle a$};
% \node at (.45,-.05) {$\scriptstyle b$};
% \end{tikzpicture}
% & = 
% \qbinomsm{a+b}{a}\:
% \begin{tikzpicture}[baseline = -1mm,scale=.6,color=\clr]
% 	\draw[-,line width=1.5pt] (0.08,-.8) to (0.08,.6);
%         \node at (.08,-1) {$\scriptstyle a+b$};
% \end{tikzpicture}, 
% \\
% \label{dotmovecrossingC}
% \begin{tikzpicture}[baseline=-1mm, color=\clr]
% 	\draw[-,thick] (0.3,-.3) to (-.3,.4);
% 	\draw[-,line width=4pt,white] (-0.3,-.3) to (.3,.4);
% 	\draw[-,thick] (-0.3,-.3) to (.3,.4);
%         \node at (-0.3,-.45) {$\scriptstyle 1$};
%         \node at (0.3,-.45) {$\scriptstyle 1$};
%      \draw(-0.15,-0.12)   \bdot;
% \end{tikzpicture}
%  & =
% \begin{tikzpicture}[baseline=-1mm, color=\clr]
%  \draw[-,thick] (-0.3,-.3) to (.3,.4);
% 	\draw[-,line width=4pt,white] (0.3,-.3) to (-.3,.4);
%   \draw[-,thick] (0.3,-.3) to (-.3,.4);
%         \node at (-0.3,-.45) {$\scriptstyle 1$};
%         \node at (0.3,-.45) {$\scriptstyle 1$};
%      \draw(0.15,0.23)   \bdot;
% \end{tikzpicture} ,  
%     \qquad
% \begin{tikzpicture}[baseline=-1mm, color=\clr]
% 	\draw[-,thick] (0.3,-.3) to (-.3,.4);
% 	\draw[-,line width=4pt,white] (-0.3,-.3) to (.3,.4);
% 	\draw[-,thick] (-0.3,-.3) to (.3,.4);
%         \node at (-0.3,-.45) {$\scriptstyle 1$};
%         \node at (0.3,-.45) {$\scriptstyle 1$};
%      \draw(-0.15,0.23)   \bdot;
% \end{tikzpicture}
%   =
% \begin{tikzpicture}[baseline=-1mm, color=\clr]
%  \draw[-,thick] (-0.3,-.3) to (.3,.4);
% 	\draw[-,line width=4pt,white] (0.3,-.3) to (-.3,.4);
%   \draw[-,thick] (0.3,-.3) to (-.3,.4);
%         \node at (-0.3,-.45) {$\scriptstyle 1$};
%         \node at (0.3,-.45) {$\scriptstyle 1$};
%      \draw(0.15,-0.12)   \bdot;
% \end{tikzpicture} .
% \end{align} 
% \end{definition}


% The $\C(q)$-linear monoidal categories $\qAWC$ and $\qAW$ are isomorphic by matching the generating objects and generating morphisms in the same symbols; cf. \cite{SSW25}.


\subsection{Cyclotomic $q$-web categories $\qW_\bfu$}
Let
\begin{equation}
\label{equ:parau}
    \bfu=(  u_1, \ldots, u_\ell)\in (\kk^*)^\ell.
\end{equation}

\begin{definition}\cite[Definition 5.2]{SSW25}
 \label{def:cycWeb}
    Suppose  $\ell\ge 1$. For any $\bfu=(u_1,\ldots, u_\ell) $ in \eqref{equ:parau}, the cyclotomic $q$-web category $\qW_\bfu$ is the quotient category of $\qAW$ by the right tensor ideal generated by $\prod_{1\le j\le \ell }g_{r}(u_j)$, for $r\ge 1$, where
    \begin{equation*}
 g_{r}(u):= 
\sum_{0\le i\le r} (-u)^i q^{-i(r-i)} \:\:
 \begin{tikzpicture}[anchorbase,scale=.7,color=\clr]
\draw[-,line width=1.5pt] (0.08,-.6) to (0.08,.5);
\node at (.08,-.8) {$\scriptstyle r$};
\draw(0.08,0) \bdot;
\draw(.7,0)node {$\scriptstyle \omega_{r-i}$};
\end{tikzpicture}.
\end{equation*}
\end{definition}

Let $ \wS_\bfu(n)$ be the path algebra of the subcategory of $\qW_\bfu$ with the object set $\Lambda_\st(n)$. Then $\wS_\bfu:=\oplus _{n\in\N}\wS_\bfu(n)$ is the path algebra of $ \qW_\bfu$. Let $\qH_\bfu$ be the full subcategory of $\qW_\bfu$ with the object set $\{(1^n)\mid n\in \N\}$. Then the path algebra of the endomorphism algebra of $(1^n)$ in $\qH_\bfu$ (or in $\qW_\bfu$)
is the cyclotomic Hecke algebra associated to $\bfu$, $\bfH_\bfu(n):=\AHe(n)\big/\langle \prod_{a=1}^\ell (X_1-u_a)\rangle$. The path algebra of  $\qH_\bfu$ is $\bfH_\bfu:=\oplus_{n\in\N}\bfH_\bfu(n) $.
In particular, 
\begin{equation}
\label{equ:cycwebhecke}
    \bfH_\bfu(n)=1_{(1^n)} \wS_\bfu(n) 1_{(1^n)}.
\end{equation}
  
For any $\lambda\in \Lambda_\st(n)$, denote by $l(\lambda)$ the length of $\la$. 
For $\lambda,\mu\in \Lambda_{\text{st}}(n)$, let 
$\Mat_{\lambda,\mu}$ be the set of all $l(\lambda)\times l(\mu)$
non-negative integer matrices $A=(a_{ij})$ with row sum vector $\la$ and column sum vector $\mu$, i.e.,  $\sum_{1\le h\le l(\mu)}a_{ih}=\lambda_i$ and $\sum_{1\le h\le l(\lambda)}a_{hj}=\mu_j$ for all $1\le i\le l(\lambda) $ and $1\le j\le l(\mu)$.  Moreover, each matrix $A\in \Mat_{\lambda,\mu}$ can be identified with a reduced chicken foot ribbon  diagram \cite[\S 4.1]{SSW25}.  Denote  
\begin{equation*}
 \PMat_{\lambda,\mu}:=\{ (A, P)\mid  A=(a_{ij})\in \Mat_{\lambda,\mu}, P=(\nu_{ij}), \nu_{ij}\in \Par_{a_{ij}} \}.   
 \end{equation*} 
 For a matrix  $ P=(\nu_{ij})_{l(\lambda)\times l(\mu)}$ of partitions, denote 
$l(P)=\max\{l(\nu_{ij})\mid 1\le i\le l(\lambda), 1\le j\le l(\mu)\}.$
We define
\begin{equation}
\label{Def-spansetof-cycqweb}    
\PMat_{\lambda,\mu}^\ell:=\{(A,P)\in\PMat_{\lambda,\mu}\mid l(P)\le \ell-1\}.
\end{equation}


For any $a\le n$ and $\lambda\in \Lambda_\st(n)$, let 
\begin{align}
\mathbf R(\lambda)_a =\Big\{(\lambda_1-a_1, \ldots, \lambda_k-a_k)\in \Lambda(n-a)~\Big |~ \sum_{j=1}^ka_j=a \Big\}.
\end{align}
Here any $\mu \in \mathbf R(\lambda)_a$ can be viewed as a strict composition by omitting zero parts.  
For any $\mu=(\lambda_1-a_1, \ldots, \lambda_k-a_k) \in \mathbf R(\lambda)_a $, let $A_\mu=(a_{ij})\in\Mat_{\lambda,\mu\cup a}$
such that 
\[ a_{i,j}=\begin{cases} \lambda_i-a_i & \text{ if $i=j$}\\
    a_i &  \text{ if $j=k+1$}\\
    0& \text{ otherwise}
\end{cases}, 
\]
where $\mu\cup a=(\lambda_1-a_1, \ldots, \lambda_k-a_k,a) $. Define a subset of $\PMat_{\lambda,\mu\cup a}^\ell$ (cf. \eqref{Def-spansetof-cycqweb}):
\begin{align}
R(\lambda,\mu)=\Big\{(A_\mu,P)\in \PMat_{\lambda,\mu\cup a}^\ell \mid P=(\nu_{ij}) \text{ with } \nu_{ij}= \emptyset \text{ unless } j=k+1
\Big\}. 
\end{align}
Denote by  $c_{\lambda,\mu} $ be the cardinality of $R(\lambda,\mu)$.
Recall $\mathbf{1}_{a,n}= \sum_{\mu\in \Lambda^a(n)}1_\mu$ from \eqref{eq:1an}.



\begin{lemma}
\label{Lem:Projective}
For any $ \eta\in \Lambda_{\text{st}}(n)$, we have an isomorphism of right $\wS_\bfu(n-a)$-modules
\[
1_\eta \wS_\bfu(n)\mathbf {1}_{a,n}\cong \bigoplus_{\lambda\in \mathbf R(\eta)_a} \big(1_\lambda \wS_\bfu(n-a)\big)^{\oplus c_{\eta,\lambda}}. 
\]
In particular, $\wS_\bfu(n)\mathbf {1}_{a,n}$ is a right projective    $\wS_\bfu(n-a)$-module.  
\end{lemma}  

\begin{proof}
  By \cite[Theorem 5.17]{SSW25},  $\Hom_{\qW_{\bfu}}(\mu,\lambda)$ has a basis given by $\PMat_{\lambda,\mu}^\ell$ \eqref{Def-spansetof-cycqweb}, for any $\lambda,\mu \in\Lambda_{\text{st}}(n)$. It follows that 
     any elements in $1_\eta \wS_\bfu(n)\mathbf{1}_{a,n}$ can be written as a linear combination of elements 
     \[ 
     (A_\lambda, P)\circ ((B,Q)\otimes 1_{(a)}) 
     \]
     with $ (A_\lambda, P)\in R(\eta,\lambda)  $ and $(B,Q)\in \PMat_{\lambda,\mu}^\ell $, for $\mu \in \Lambda_{\text{st}}(n-a)$.
     Hence 
     \[
     1_\eta \wS_\bfu(n)\mathbf{1}_{a,n}=\bigoplus_{\lambda\in \mathbf R(\eta)_a, (A_\lambda, P)\in R(\eta,\lambda) } (A_\lambda, P)(\wS_\bfu(n-a)\otimes 1_{(a)}).
     \]
     Now the result follows since 
     \[ (A_\lambda, P)(\wS_\bfu(n-a)\otimes 1_{(a)})  \cong 1_\lambda \wS_\bfu(n-a) \]
     as a right $\wS_\bfu(n-a)$-module, which is projective. 
 \end{proof}



\subsection{Restriction functors and induction functors for $\qW_\bfu$}

Recall $\iota_{a,n}$  in \eqref{equ:embedding} for $\qAW$. Similarly, by the basis theorem of $\qW_\bfu$ in \cite[Theorem 5.17]{SSW25},  there is an obvious embedding from $\wS_\bfu(n-a) \rightarrow \mathbf{1}_{a,n}\wS_\bfu(n)\mathbf{1}_{a,n}$ for any $a\le n$. Then   
we introduce the induction functors and restriction functors 
\begin{align}  \label{eq:IndRes}
\begin{split}
\Ind_{n-a}^{n}:= \wS_\bfu(n)\mathbf{1}_{a,n}\otimes _{\wS_\bfu(n-a)} ? & : \wS_\bfu(n-a)\modf \longrightarrow \wS_\bfu(n)\modf. 
\\
\text{Res}^{n}_{n-a}=\mathbf{1}_{a,n}\wS_\bfu(n)\otimes _{\wS_\bfu(n)}? &: \wS_\bfu(n)\modf\longrightarrow \wS_\bfu(n-a)\modf.
\end{split}
\end{align}

\begin{lemma}
 The functors $\Ind^n_{n-a}$ and $\Res^{n}_{n-a}$ are exact.
\end{lemma}

\begin{proof}
  The exactness of $\Ind^n_{n-a}$ follows from Lemma  \ref{Lem:Projective} since    $\wS_\bfu(n)\mathbf{1}_{a,n}$ is a right projective $\wS_\bfu(n-a)$-module. The restriction functor is clearly exact.
\end{proof}
